\documentclass[11pt,a4paper]{article}
\usepackage[utf8]{inputenc}
\usepackage{amsmath,amssymb,amsthm}
\usepackage{listings}
\usepackage[margin=2.5cm]{geometry}
\usepackage{xcolor}

\newtheorem{theorem}{Theorem}

\lstset{
  language=C++,
  basicstyle=\ttfamily\small,
  keywordstyle=\color{blue}\bfseries,
  commentstyle=\color{green!50!black},
  stringstyle=\color{red!70!black},
  numbers=left,
  numberstyle=\tiny\color{gray},
  breaklines=true,
  frame=single,
  tabsize=4,
  showstringspaces=false
}

\title{IOI 1990 -- Problem 2: Map Labelling}
\author{Solution}
\date{}

\begin{document}
\maketitle

\section{Problem Statement}

Given $n$ points on a 2D map, each requiring a rectangular label of specified
dimensions, place each label adjacent to its point (in one of 2 candidate
positions) such that no two labels overlap. Determine whether a valid
labelling exists, and if so, output one.

\section{Solution}

\subsection{2-SAT Formulation}

When each point has exactly two candidate positions, this is a classic
\textbf{2-SAT} problem:

\begin{itemize}
  \item For each point $i$, define a boolean variable $x_i$:
    $x_i = \texttt{true}$ means position~0, $x_i = \texttt{false}$ means
    position~1.
  \item For each pair $(i, j)$ with $i < j$, check all 4 combinations
    of positions. If positions $(p_i, p_j)$ cause an overlap, add the
    clause $(\neg x_i^{p_i} \lor \neg x_j^{p_j})$, which forbids both
    choices simultaneously.
\end{itemize}

\subsection{Solving 2-SAT via Kosaraju's Algorithm}

\begin{enumerate}
  \item \textbf{Build the implication graph.} For each clause
    $(\ell_1 \lor \ell_2)$, add directed edges $\neg\ell_1 \to \ell_2$
    and $\neg\ell_2 \to \ell_1$.
  \item \textbf{First DFS pass} on the original graph to compute a
    topological finishing order.
  \item \textbf{Second DFS pass} on the reverse graph, processing vertices
    in reverse finishing order, to identify strongly connected components
    (SCCs).
  \item \textbf{Check satisfiability:} if $x_i$ and $\neg x_i$ belong to
    the same SCC for any~$i$, the instance is unsatisfiable.
  \item \textbf{Extract assignment:} set $x_i = \texttt{true}$ if $x_i$'s
    SCC has a higher topological index than $\neg x_i$'s SCC.
\end{enumerate}

\begin{theorem}
A 2-SAT formula is satisfiable if and only if no variable and its negation
belong to the same SCC in the implication graph. When satisfiable, the SCC
ordering yields a valid assignment.
\end{theorem}

\subsection{Overlap Detection}

Two axis-aligned rectangles with corners $(lx_i, ly_i)$ and
$(lx_i + w_i, ly_i + h_i)$ overlap if and only if:
\[
  lx_i < lx_j + w_j \;\land\; lx_j < lx_i + w_i
  \;\land\; ly_i < ly_j + h_j \;\land\; ly_j < ly_i + h_i.
\]

\section{Complexity Analysis}

\begin{itemize}
  \item \textbf{Time:} $O(n^2)$ for generating clauses (checking all pairs),
    plus $O(n + m)$ for the two DFS passes where $m = O(n^2)$ is the number
    of implication edges. Total: $O(n^2)$.
  \item \textbf{Space:} $O(n^2)$ for the implication graph.
\end{itemize}

\section{Implementation}

\begin{lstlisting}
#include <iostream>
#include <vector>
#include <algorithm>
using namespace std;

const int MAXN = 2005;

struct TwoSat {
    int n;
    vector<int> adj[2 * MAXN], radj[2 * MAXN];
    int comp[2 * MAXN];
    bool visited[2 * MAXN];
    vector<int> topo;

    void init(int _n) {
        n = _n;
        for (int i = 0; i < 2 * n; i++) {
            adj[i].clear();
            radj[i].clear();
        }
    }

    // Add clause (a OR b). Literal for variable i:
    // true = 2*i, false = 2*i+1.
    void addClause(int a, int b) {
        adj[a ^ 1].push_back(b);
        adj[b ^ 1].push_back(a);
        radj[b].push_back(a ^ 1);
        radj[a].push_back(b ^ 1);
    }

    void dfs1(int u) {
        visited[u] = true;
        for (int v : adj[u])
            if (!visited[v]) dfs1(v);
        topo.push_back(u);
    }

    void dfs2(int u, int c) {
        comp[u] = c;
        for (int v : radj[u])
            if (comp[v] == -1) dfs2(v, c);
    }

    bool solve(vector<bool>& result) {
        fill(visited, visited + 2 * n, false);
        topo.clear();
        for (int i = 0; i < 2 * n; i++)
            if (!visited[i]) dfs1(i);
        fill(comp, comp + 2 * n, -1);
        int c = 0;
        for (int i = 2 * n - 1; i >= 0; i--) {
            int u = topo[i];
            if (comp[u] == -1) dfs2(u, c++);
        }
        result.resize(n);
        for (int i = 0; i < n; i++) {
            if (comp[2 * i] == comp[2 * i + 1])
                return false;
            result[i] = (comp[2 * i] > comp[2 * i + 1]);
        }
        return true;
    }
};

struct Point {
    double x, y;
    double w, h;
    // Position 0: label to the right
    // Position 1: label to the left
    double lx(int pos) const {
        return pos == 0 ? x : x - w;
    }
    double ly(int /*pos*/) const {
        return y;
    }
};

bool overlaps(const Point& a, int pa, const Point& b, int pb) {
    double ax = a.lx(pa), ay = a.ly(pa);
    double bx = b.lx(pb), by = b.ly(pb);
    return ax < bx + b.w && bx < ax + a.w &&
           ay < by + b.h && by < ay + a.h;
}

int main() {
    int n;
    cin >> n;

    vector<Point> pts(n);
    for (int i = 0; i < n; i++)
        cin >> pts[i].x >> pts[i].y >> pts[i].w >> pts[i].h;

    TwoSat sat;
    sat.init(n);

    for (int i = 0; i < n; i++) {
        for (int j = i + 1; j < n; j++) {
            for (int pi = 0; pi < 2; pi++) {
                for (int pj = 0; pj < 2; pj++) {
                    if (overlaps(pts[i], pi, pts[j], pj)) {
                        // Forbid (x_i = pi) AND (x_j = pj).
                        // Literal for "x_i = pi": 2*i if pi=0,
                        //                         2*i+1 if pi=1.
                        // Negation: 2*i+1 if pi=0, 2*i if pi=1.
                        int li = (pi == 0) ? (2*i+1) : (2*i);
                        int lj = (pj == 0) ? (2*j+1) : (2*j);
                        sat.addClause(li, lj);
                    }
                }
            }
        }
    }

    vector<bool> result;
    if (sat.solve(result)) {
        cout << "YES" << endl;
        for (int i = 0; i < n; i++) {
            int pos = result[i] ? 0 : 1;
            cout << "Point " << i+1 << ": position " << pos
                 << " (" << pts[i].lx(pos) << ", "
                 << pts[i].ly(pos) << ")" << endl;
        }
    } else {
        cout << "NO" << endl;
    }

    return 0;
}
\end{lstlisting}

\section{Notes}

\begin{itemize}
  \item The 2-SAT approach handles the case of exactly 2 candidate positions
    per point. For 4 positions per point, the problem becomes NP-hard in
    general.
  \item Kosaraju's algorithm requires two full DFS passes, which is simple
    to implement. Tarjan's single-pass SCC algorithm is an alternative.
  \item For floating-point coordinates, exact overlap detection may require
    care with precision. Using integer coordinates and strict inequalities
    avoids this issue in most contest settings.
\end{itemize}

\end{document}
